\chapter{Marco Metodológico}\label{metodologia}

A continuación, se detalla la metodología de trabajo seguida para el desarrollo del proyecto, estructurada en fases semanales según el cronograma establecido. Cada fase se enfocó en el cumplimiento de objetivos específicos para asegurar una construcción modular y robusta del sistema de control.

\begin{enumerate}[label=\textbf{Semana \arabic*:}, leftmargin=*]

    \item \textbf{Revisión bibliográfica y configuración del entorno (25/08 al 29/08):}
    Se realizó una investigación inicial donde se estudiaron las hojas de datos de los componentes clave: el microcontrolador ESP32, el driver de motor de paso DM542T y el encoder incremental. Se configuró el entorno de desarrollo Visual Studio Code con la extensión PlatformIO y el framework ESP-IDF, realizando pruebas de compilación y carga de código para verificar la comunicación con la placa de desarrollo. Como añadidura se investigaron los principios de control para un péndulo invertido.

    \item \textbf{Control básico de periféricos y calibración (01/09 al 05/09):}
    Se desarrollaron módulos de software independientes (librerías) para el control de bajo nivel. Se implementó una librería para el motor de paso basada en \textbf{ledc} capaz de generar trenes de pulsos a una frecuencia variable. Paralelamente, se creó una librería para el encoder incremental, utilizando el periférico de hardware PCNT del ESP32 en modo cuadratura para una lectura precisa del ángulo.

    \item \textbf{Implementación de la lógica de control (08/09 al 12/09):}
    Se diseñó la arquitectura del sistema de control principal. Se implementó inicialmente un seguidor, es decir, el motor debía girar tanto como se giraba el encoder, de forma sincronizada, utilizando un controlador proporcional, estableciendo un conjunto de tareas en FreeRTOS para el bucle de control. Se realizaron pruebas en lazo abierto (comandos manuales) a travéz de uart con el módulo de software "uart\_echo" para verificar la respuesta del actuador y, posteriormente, se cerró el lazo utilizando la realimentación del encoder. Se inició el proceso de sintonización empírica de la ganancia proporcional.

    \item \textbf{Diseño de la interfaz de visualización (15/09 al 19/09):}
    Se integró el hardware de la pantalla LCD mediante comunicación THT de forma paralela. Se desarrolló un módulo de software "lcd\_controller" para abstraer las operaciones de bajo nivel.
    Por otra parte se incorporó el sistema de control a un riel de pruebas, el cual llevó a la creación de un módulo de software, llamado "button\_handler", para gestionar la lectura de botones físicos que permitieran establecer los márgenes físico de parada por protección.

    \item \textbf{Integración del sistema (22/09 al 26/09):}
    Se unificaron todos los módulos desarrollados (controlador, motor, encoder, LCD, botones) en una única aplicación concurrente,. Se ajustaron las prioridades de las tareas para garantizar que el bucle de control en tiempo real no fuera afectado por las tareas de menor prioridad, como la actualización de la pantalla o la lectura de comandos UART.

    \item \textbf{Optimización y pruebas de robustez (29/09 al 03/10):}
    Se realizó un proceso iterativo de sintonización fina de las ganancias del controlador para mejorar la estabilidad y la respuesta ante perturbaciones. Se implementaron optimizaciones en el código, en cuanto a errores del sistema corregidos, y se realizaron pruebas de robustez para verificar el comportamiento del sistema en los límites de operación y la fiabilidad de los mecanismos de seguridad. Por otro lado se realizó el cambio de comunicación paralela a I2C para la pantalla LCD, liberando pines GPIO y mejorando la eficiencia del sistema.

    \item \textbf{Diseño de módulo de pruebas (06/10 al 10/10):}
    Se diseñó un módulo de pruebas que tenga como función principal el control del driver del motor a través de un esp32, sin embargo también se puede usar en conjunto con una pantalla LCD, un encoder y switches externos. Se incorporó una interfaz de usuario con múltiples vistas, gestionada por una tarea de baja prioridad, para mostrar en tiempo real el estado del sistema, el ángulo del péndulo, y las constantes del controlador, facilitando el diagnóstico y la sintonización.

    \item \textbf{Documentación del código y manual de usuario (13/10 al 17/10):}
    Se llevó a cabo una revisión completa del código fuente para añadir comentarios explicativos en formato estándar. Se documentó la arquitectura general del software, el propósito de cada módulo (librería), la función de cada tarea de FreeRTOS y la interfaz pública de las funciones más importantes para facilitar su comprensión y futuro mantenimiento.
    Se recopilaron los datos de rendimiento obtenidos durante las pruebas por separado del sistema y se obtuvieron las gráficas de pulsos para visualizar el comportamiento del sistema. Se preparó el material audiovisual para la presentación de resultados y se redactó el informe técnico final del proyecto, detallando el diseño, la implementación, las pruebas realizadas y las conclusiones. Por otra parte, se elaboró un manual de usuario que incluye instrucciones para la instalación, configuración y operación del sistema, así como recomendaciones para el mantenimiento y solución de problemas comunes.

\end{enumerate}